\chapter{Exploration}
\label{ch:exploration}

An agent learns only about what it tries. The bandits of \cref{ch:bandits}
showed the dilemma in its simplest form and the main answers: random
exploration ($\varepsilon$-greedy), optimism (UCB) and posterior sampling
(Thompson). In an MDP the dilemma is harder, because information is
delayed: reaching an informative state may take many deliberate steps, and
random actions rarely string them together. This chapter shows why, and
presents the methods that scale the bandit ideas to deep RL.

\section{Why dithering fails}

$\varepsilon$-greedy, Boltzmann exploration and the entropy bonus of PPO
are \emph{dithering} strategies: they perturb each action independently.
The following example, a ``combination lock'', shows that they can need a
time exponential in the size of the problem.

\begin{proposition}{Random exploration in a chain}{chain}
  Consider states $0, 1, \dots, n$, starting in $0$, where action
  \emph{right} moves from $k$ to $k + 1$ and action \emph{left} returns to
  $0$; the only reward is at state $n$. Under the uniform random policy,
  the expected number of steps to reach $n$ from $0$ is $2^{n+1} - 2$.
\end{proposition}

\begin{proof}
  Let $T_k$ be the expected number of steps to reach $n$ from $k$, so that
  $T_n = 0$ and $T_k = 1 + \frac{1}{2} T_{k+1} + \frac{1}{2} T_0$ for $k <
  n$. We show by induction that $T_k = T_0 - (2^{k+1} - 2)$. It holds for
  $k = 0$, and if it holds for $k$, then $T_{k+1} = 2 T_k - 2 - T_0 = T_0 -
  2^{k+2} + 2$. With $k = n$, $0 = T_0 - 2^{n+1} + 2$.
\end{proof}

A greedy agent with optimistic initial values solves the same chain in
$O(n^2)$ steps: every unvisited action looks attractive, so the agent
tries the actions of the states it reaches systematically. This is the
principle of the provably efficient algorithms.

\section{Optimism in the face of uncertainty}

\paragraph{Provably efficient exploration.} R-max (Brafman and
Tennenholtz, 2002) treats every state-action pair visited fewer than $m$
times as leading to a fictitious state of maximal reward $R_{\max}$, and
plans in the resulting model; either the plan collects near-optimal reward,
or it visits an under-explored pair, which can happen only finitely often.
R-max and its successors (E$^3$, MBIE, UCRL2) find a near-optimal policy
with a number of mistakes polynomial in the numbers of states and actions,
in $1 / (1 - \gamma)$ and in the accuracy: the MDP analogue of the
logarithmic regret of UCB.

\paragraph{Exploration bonuses.} A simpler form of optimism adds a bonus
to the reward,
\begin{equation}
  \label{eq:bonus}
  \tilde r(s, a) = r(s, a) + \frac{\beta}{\sqrt{N(s, a)}},
\end{equation}
where $N(s, a)$ is the number of visits (MBIE-EB, Strehl and Littman,
2008). The square root has the same origin as in UCB1 (\cref{def:ucb1}):
the confidence interval of a mean estimated from $N$ samples has a width
of order $1 / \sqrt{N}$. The bonus is propagated by the Bellman equations,
so the agent is attracted to states from which under-explored pairs can be
reached, not only to these pairs themselves.

\section{Counts in large state spaces}

With images as states, every state is new, and $N(s) \in \{0, 1\}$. Deep
RL methods replace counts by quantities that generalise.

\paragraph{Pseudo-counts.} Bellemare et al.\ (2016) derive a count from a
density model $\rho_n(s)$ trained on the states seen so far. Let $\rho'_n(s)$
be the density the model would give to $s$ after one more training step on
$s$. If the model were an empirical frequency, $\rho_n(s) = N / n$ and
$\rho'_n(s) = (N + 1) / (n + 1)$. Solving these two equations for $N$
defines the \emph{pseudo-count}
\begin{equation}
  \hat N_n(s) = \frac{\rho_n(s) \left( 1 - \rho'_n(s) \right)}{\rho'_n(s) -
  \rho_n(s)},
\end{equation}
used in the bonus \eqref{eq:bonus}. (From $\rho = N/n$ and $\rho' = (N +
1)/(n + 1)$: $n = N / \rho$, so $\rho' (N/\rho + 1) = N + 1$, which gives
$N (\rho' - \rho) = \rho (1 - \rho')$.) With a pixel density model, the
method made the first real progress on \emph{Montezuma's Revenge}, an
Atari game where rewards are rare and DQN scored zero.

\paragraph{Prediction error as novelty.} Random network distillation
(RND, Burda et al., 2019) trains a network $\hat f_\psi$ to predict the
output of a fixed, randomly initialised network $f$, on the states visited.
The error $\norm{\hat f_\psi(s) - f(s)}^2$ is small on states similar to
those seen often, and large on new states: it is used as the bonus. The
target is deterministic, so the error does not stay large on states whose
future is merely random, a failure of the earlier \emph{curiosity}
methods, which use the error of a forward model of the dynamics. The
intrinsic curiosity module (ICM, Pathak et al., 2017) avoids part of the
problem by predicting in a feature space learnt to predict the agent's own
actions (an inverse model), which ignores what the agent cannot influence:
a television showing random images is not interesting to it.

\begin{lstlisting}[caption={Random network distillation bonus (PyTorch)}]
target = make_net().requires_grad_(False)   # fixed random network
predictor = make_net()
opt = torch.optim.Adam(predictor.parameters(), lr=1e-4)

def bonus(states):                          # also trains the predictor
    err = (predictor(states) - target(states)).pow(2).sum(1)
    opt.zero_grad(); err.mean().backward(); opt.step()
    return err.detach()                      # intrinsic reward
\end{lstlisting}

In practice the intrinsic rewards are normalised, combined with the
extrinsic rewards with their own discount and value head, and used with
PPO.

\paragraph{Remembering where to return.} Go-Explore (Ecoffet et al.,
2021) observed that novelty bonuses fade before the agent has learnt to
return to the frontier of the explored region. It keeps an archive of
visited ``cells'' (downscaled images), returns to a promising cell
(deterministically in a simulator, or with a goal-conditioned policy), and
explores from there; the trajectories found are then made robust by
imitation. It solved \emph{Montezuma's Revenge} and \emph{Pitfall}.

\section{Posterior sampling and deep exploration}

Thompson sampling also extends to MDPs. Posterior sampling for RL (PSRL,
Osband et al., 2013) keeps a posterior over MDPs, samples one at the start
of each episode, and follows its optimal policy for the whole episode.
Committing to one hypothesis for an episode produces \emph{deep}
exploration: if the sampled MDP is optimistic about a distant state, the
agent goes there deliberately, unlike dithering. PSRL has a regret of
order $\tilde O(H S \sqrt{A T})$ for episodes of length $H$.

Bootstrapped DQN (Osband et al., 2016) approximates the posterior with an
ensemble of $K$ Q-heads on a shared torso, each trained on its own
bootstrap sample of the data (each transition is used by each head with
probability $1/2$). At each episode one head is drawn at random and
followed greedily. NoisyNets (\cref{sec:rainbow}) are a cheaper relative:
the noise of the weights is sampled once for many steps, which also gives
temporally consistent exploration.

\begin{pacman}{Exploration in Pacman}{exploration-pacman}
  Exploration is not the bottleneck in Pacman: the cookies are a dense
  reward, and the random ghosts produce a wide variety of situations by
  themselves. $\varepsilon$-greedy (DQN) and the entropy bonus with
  $c_e = 0.03$ (PPO) suffice. Where exploration matters is in the rare,
  dangerous situations (being trapped between two ghosts), which the agent
  meets too seldom to learn well: this is why the look-ahead search of
  \cref{ch:model-based} helps most there. A maze with a single, distant
  reward and no cookies would be a hard exploration problem.
\end{pacman}

\section{Summary}

\begin{itemize}[nosep]
  \item Dithering ($\varepsilon$-greedy, entropy bonuses) explores in time
    exponential in the depth of the reward, even in a chain.
  \item Optimism (R-max, bonuses $\beta / \sqrt{N}$) gives polynomial
    guarantees; in large spaces, counts are replaced by pseudo-counts from
    density models or by prediction errors (RND, ICM).
  \item Posterior sampling (PSRL, bootstrapped DQN) explores deeply by
    committing to a sampled hypothesis for a whole episode.
\end{itemize}
